\section{Results}

The main result is a reversal under cost matching. In the initial five-seed reliability check, LV-VoI scale3 appeared plausible because it achieved a success rate close to a higher-cost random baseline. That interpretation changed once the random family was swept around the method's realized human cost. As shown in Table~\ref{tab:registry_costmatched_claims}, random b350 achieved 439/500 repeated successes (87.8\%, 95\% CI [84.6, 90.4]) at 177.0 best-step human steps, whereas LV-VoI scale3 achieved 416/500 repeated successes (83.2\%, 95\% CI [79.7, 86.2]) at 202.0 best-step human steps. The learned trigger is therefore not on the observed Lift success-cost frontier, so the current evidence does not support a claim that it allocates human attention more efficiently than random intervention.

\begin{figure*}[t]
    \centering
    \includegraphics[width=0.95\textwidth]{fig1_methodology_protocol_r056.pdf}
    \caption{Methodology-first overview of the diagnostic evaluation route. The figure states the gate sequence for HIL-RL human-attention allocation claims, then shows the current Lift verdicts and derived stop-rule metrics from registered R021/R023/R024 evidence. It is a protocol visualization rather than a new empirical result.}
    \label{fig:protocol_hero}
\end{figure*}

\begin{table}[t]
\centering
\caption{Registry-generated cost-matched Lift claims. All numeric entries are drawn from the evidence registry.}
\label{tab:registry_costmatched_claims}
\begin{tabular}{llcccc}
\toprule
Strategy & Run & Success & 95\% CI & Human steps & $\Delta$ \\
\midrule
No intervention & R020 & 400/500 = 80.0\% & [76.3, 83.3] & 0.0 & -3.2 pp \\
\textbf{Random b350} & R021 & 439/500 = 87.8\% & [84.6, 90.4] & 177.0 & +4.6 pp \\
LV-VoI scale3 & R021 & 416/500 = 83.2\% & [79.7, 86.2] & 202.0 & +0.0 pp \\
Random b450 & R021 & 426/500 = 85.2\% & [81.8, 88.0] & 250.0 & +2.0 pp \\
Random b600 & R021 & 426/500 = 85.2\% & [81.8, 88.0] & 269.2 & +2.0 pp \\
\bottomrule
\end{tabular}
\end{table}


The reversal is not an artifact of including only one random budget. Random b450 and random b600 provide higher-cost random-family points, while random b350 provides the lower-cost point that invalidates the original efficiency interpretation. This matters because LV-VoI scale3 would look more favorable if it were compared only with the original higher-cost random condition. The result therefore supports the first protocol gate: a human-attention allocation claim should be tested against a local random success-cost frontier, not against a single random baseline selected before the realized method cost is known.

Two lightweight trigger repairs remain negative under the same discipline. The minimum-disagreement repair reached 226/300 repeated successes (75.3\%, 95\% CI [70.2, 79.9]) at 211.7 best-step human steps. The score-floor repair reached 233/300 repeated successes (77.7\%, 95\% CI [72.6, 82.0]) at 253.3 best-step human steps. Both are below the same-seed random b350 reference, which reached 259/300 repeated successes (86.3\%, 95\% CI [82.0, 89.8]) at 95.0 best-step human steps. Table~\ref{tab:registry_trigger_repairs} is generated from the audited registry and is the preferred numeric source for this repair comparison.

\begin{table}[t]
\centering
\caption{Registry-generated negative trigger-repair claims on Lift seeds 0--2.}
\label{tab:registry_trigger_repairs}
\begin{tabular}{llcccc}
\toprule
Strategy & Run & Success & 95\% CI & Human steps & $\Delta$ \\
\midrule
\textbf{Random b350} & R024 & 259/300 = 86.3\% & [82.0, 89.8] & 95.0 & +0.0 pp \\
Min-disagree LV-VoI & R022 & 226/300 = 75.3\% & [70.2, 79.9] & 211.7 & -11.0 pp \\
Score-floor LV-VoI & R024 & 233/300 = 77.7\% & [72.6, 82.0] & 253.3 & -8.7 pp \\
\bottomrule
\end{tabular}
\end{table}


\begin{figure*}[t]
    \centering
    \includegraphics[width=0.95\textwidth]{fig_attention_allocation_diagnostics_r054.pdf}
    \caption{Attention-allocation diagnostics for the current negative trigger result. The panels combine the R021 cost-matched reversal, the R024 repair stop gate, and R023/R024 intervention-start traces to show when budget is spent, how much budget has already been used, where the gripper is relative to the cube, and whether LV-VoI scores or failure probabilities saturate. The figure is diagnostic: it visualizes the registered R021/R023/R024 evidence and does not introduce a new positive method claim.}
    \label{fig:attention_allocation_diagnostics_r054}
\end{figure*}

The trace results sharpen the mechanism claim. Figure~\ref{fig:attention_allocation_diagnostics_r054} summarizes the trace-level evidence without treating it as a new success metric. LV-VoI does not fail merely because it requests interventions far from the cube. In the R023 traces, LV-VoI starts were closer to the cube than random on average, with mean gripper-to-cube norm 0.217 for LV-VoI versus 0.285 for random b350, and mean gripper-to-cube planar distance 0.152 versus 0.235. However, LV-VoI used 96 intervention starts across seeds 0--2 compared with 55 starts for random b350. The trigger therefore appears to spend more budget despite querying in geometrically plausible states, which points away from a simple distance-threshold repair.

The score-floor repair confirms this interpretation. R024 mechanically removed low-score starts after step 4000: the trace table records zero post-floor starts below the 0.05 threshold. Nevertheless, the repair still produced 94 intervention starts, almost the same as the original LV-VoI count of 96 and far above the random b350 count of 55. Its median accepted score was only 0.067, indicating that the floor was too weak to create random-level selectivity. Together, the repair and trace results support a narrower conclusion: the current trigger suffers from budget-control and score-timing mismatch, and any new trigger should be evaluated through the same cost-matched and trace-level protocol before being expanded.

Stack provides boundary evidence rather than a positive generalization result. With a strong matched behavioral-cloning initialization, the no-online-human Stack baseline achieved 131/180 repeated successes (72.8\%) at zero online human cost, while both random intervention and Stack-tuned LV-VoI achieved 107/180 repeated successes (59.4\%) at substantial human cost. These results do not support a broad robotics-transfer claim for the current trigger. They instead reinforce the paper's conservative framing: the Lift evidence motivates the diagnostic protocol, and Stack shows that adding online interventions can be harmful when the task and update mechanism are not aligned.
